\chapter{Pricing Models}

\section{How Competitors Price}

Competitor pricing falls into two camps: outright purchase and
Robots-as-a-Service. Purchase prices range from about USD 30,000--45,000 for a
MiR base \cite{mir250} and around EUR 50,000 for a Robotnik RB-KAIROS+
\cite{robotnik_kairos}, up to USD 75,000+ for a Boston Dynamics Spot and USD
120,000--320,000 for high-end security platforms
\cite{standardbots_amr,mordor_security}. Full warehouse deployments of 50--100
robots run into the millions. On the service side, players such as Locus and
Cobalt sell subscriptions (Cobalt around USD 3,500--6,500 per month), which
shift the customer from a large capital outlay to an operating expense
\cite{builtin_raas,coherent_security}. A structured summary is given in
Table~\ref{tab:pricing}.

\section{Why Outright Purchase Fails Locally}

The Tunisian problem is not the sticker price in isolation but the
robot-cost-to-labor-cost ratio. Against low local wages, an outright purchase
has a much longer payback than it does in Europe, and the upfront cash plus
foreign-currency exposure on imported hardware makes a large capital purchase
hard for a Tunisian SME to approve. Outright sale is therefore the wrong default
model here.

\section{Recommended Pricing for the Local Market}

\begin{itemize}
    \item \textbf{Robots-as-a-Service / leasing} as the default, so the customer
          pays from an operating budget instead of committing capital
          \cite{builtin_raas}.
    \item \textbf{Usage-based pricing} (per inspection round, per hour, per
          delivery) where the value is measurable, aligning cost with benefit.
    \item \textbf{Shared deployments} between several smaller companies in the
          same zone, spreading the cost of one robot across more than one
          buyer -- a model that fits Tunisia's SME structure and that no global
          competitor offers.
    \item \textbf{Pilot-first}: low-cost or free initial pilots in universities
          and private clinics to generate references quickly, on the
          understanding that the first reference is worth more than the first
          margin.
\end{itemize}

\begingroup
\small
\setlength{\tabcolsep}{4pt}
\begin{longtable}{|p{2.7cm}|p{2.2cm}|p{1.9cm}|p{2.5cm}|p{1.8cm}|p{1.3cm}|}
\caption{Pricing models across competitors, with the recommended local model.}\label{tab:pricing}\\
\hline
\textbf{Company / Product} & \textbf{Pricing Model} & \textbf{Estimated Price} & \textbf{Included Services} & \textbf{Notes} & \textbf{Source} \\
\hline
\endfirsthead
\hline
\textbf{Company / Product} & \textbf{Pricing Model} & \textbf{Estimated Price} & \textbf{Included Services} & \textbf{Notes} & \textbf{Source} \\
\hline
\endhead
\hline
\endfoot
MiR250/600 & Purchase (hardware) & USD 30--45k & Base + interchangeable top modules & Proprietary software lock-in & \cite{mir250} \\
\hline
Robotnik RB-KAIROS+ & Purchase & From EUR 50k & ROS2 base + plug-and-play arm & High entry price & \cite{robotnik_kairos} \\
\hline
Boston Dynamics Spot & Purchase & USD 75k+ & Legged base + sensor mounts & High cost and training overhead & \cite{standardbots_amr} \\
\hline
SMP Robotics & Purchase & USD 120--320k & Outdoor security patrol & Extremely high cost & \cite{mordor_security} \\
\hline
Cobalt Robotics & Subscription (RaaS) & USD 3.5--6.5k/mo & Indoor patrol + monitoring & Capex-light; security-only & \cite{coherent_security} \\
\hline
Locus Robotics & Subscription (RaaS) & Per-bot subscription & Collaborative picking fleet & Capex-light; single workflow & \cite{builtin_raas} \\
\hline
Large warehouse fleet & Purchase / project & USD 2--4M (50--100 bots) & Full goods-to-person system + integration & Out of scale for Tunisian SMEs & \cite{valuates_agv_amr} \\
\hline
\ProjectName{} (recommended) & RaaS / usage-based / shared / pilot & TBD (local) & Robot + local support, maintenance, training & Designed for the local cost-to-labor ratio & TBD (local) \\
\hline
\end{longtable}
\endgroup

